\documentclass[12pt,a4paper]{article}

\usepackage[utf8]{inputenc}
\usepackage[T1]{fontenc}
\usepackage[spanish]{babel}
\usepackage{geometry}
\usepackage{hyperref}
\usepackage{listings}
\usepackage{xcolor}
\usepackage{graphicx}
\usepackage{tabularx}
\usepackage{enumitem}
\usepackage{fancyhdr}

\geometry{margin=2.5cm}

\definecolor{codebg}{rgb}{0.95,0.95,0.95}
\definecolor{codeblue}{rgb}{0.2,0.2,0.6}
\definecolor{codegreen}{rgb}{0.1,0.5,0.1}

\lstset{
    backgroundcolor=\color{codebg},
    basicstyle=\ttfamily\small,
    keywordstyle=\color{codeblue},
    commentstyle=\color{codegreen},
    frame=single,
    breaklines=true,
    tabsize=2,
    captionpos=b
}

\hypersetup{
    colorlinks=true,
    linkcolor=blue,
    urlcolor=blue
}

\pagestyle{fancy}
\lhead{Vetrinaria}
\rhead{Sprints 0--10}
\cfoot{\thepage}

\title{\vspace{-2cm}\textbf{Vetrinaria}\\[0.3cm]
\large Sistema de Gestión para Clínicas Veterinarias\\[0.5cm]
\textbf{Documentación Técnica — Sprints 0 a 10}}
\author{}
\date{Mayo 2026}

\begin{document}
\maketitle
\thispagestyle{empty}
\newpage

\tableofcontents
\newpage

\section{Resumen Ejecutivo}

Este documento detalla el trabajo realizado durante los sprints del
proyecto \textbf{Vetrinaria}, un sistema integral para la gestión de clínicas
veterinarias.

\begin{itemize}[leftmargin=*]
    \item \textbf{Sprint 0:} Configuración del monorepo, base de datos, Docker,
          componentes UI y diseño responsivo.
    \item \textbf{Sprint 1:} Sistema de autenticación con NextAuth v5, páginas de
          login/registro y middleware de protección por roles.
    \item \textbf{Sprint 2:} CRUD completo de Clientes (dueños de mascotas).
    \item \textbf{Sprint 3:} CRUD completo de Pacientes (mascotas vinculadas a clientes).
    \item \textbf{Sprint 4:} Agenda de Citas --- vista diaria, CRUD completo con
          transiciones de estado y selectores dependientes.
    \item \textbf{Sprint 5:} Historial Clínico --- registros SOAP, constantes
          vitales y vinculación con pacientes/citas.
    \item \textbf{Sprint 6:} Facturación --- CRUD de productos, proveedores, facturas
          con IVA 16\% y generación de número secuencial.
    \item \textbf{Sprint 7:} Dashboard avanzado con métricas reales desde la base
          de datos (citas, clientes, pacientes, ingresos).
    \item \textbf{Sprint 8:} Notificaciones --- tabla, migración, servidor de
          notificaciones, campana en Navbar con contador en vivo y página dedicada.
    \item \textbf{Sprint 9:} Backups --- creación, listado, restauración y
          eliminación de copias de seguridad via pg\_dump desde la interfaz web.
    \item \textbf{Sprint 10:} Multi-tenant --- arquitectura multi-clínica con
          tabla \texttt{clinics}, migración de esquemas, filtrado por clínica
          en consultas críticas y gestión desde configuración.
\end{itemize}

\section{Sprint 0 — Setup del Proyecto}

\subsection{Arquitectura del Monorepo}

El proyecto utiliza \textbf{Turborepo} con \textbf{pnpm} como gestor de paquetes,
organizado en la siguiente estructura:

\begin{lstlisting}[language=bash, caption=Estructura del monorepo]
vetrinaria/
+-- apps/
|   +-- web/                    # Next.js 15 App Router
|       +-- src/
|           +-- app/            # Rutas y layouts
|           +-- components/     # Componentes React
|           +-- lib/            # Utilidades y config
+-- packages/
|   +-- db/                     # Drizzle ORM + esquemas
|   +-- shared/                 # Tipos y constantes
|   +-- config/                 # ESLint, tsconfig base
+-- docker-compose.yml          # Servicios PostgreSQL, MinIO, Redis
+-- turbo.json                  # Configuración de Turborepo
\end{lstlisting}

Las dependencias se gestionan con \texttt{pnpm-workspace.yaml}, que define los
workspaces del monorepo. Turborepo permite ejecutar tareas (\texttt{build},
\texttt{lint}, \texttt{typecheck}) en paralelo con caching inteligente.

\subsection{Configuración de Next.js 15 y TypeScript}

Se configuró \textbf{Next.js 15} con el \textbf{App Router} y \textbf{TypeScript 5.8},
utilizando configuraciones base compartidas desde \texttt{packages/config}:

\begin{itemize}[leftmargin=*]
    \item \texttt{tsconfig.base.json}: Configuración base con ES2022, strict mode,
          y bundler module resolution.
    \item \texttt{tsconfig.nextjs.json}: Extiende la base con jsx:preserve y
          el plugin de Next.js.
\end{itemize}

El build de Next.js 15 compila exitosamente en aproximadamente 30 segundos
en el entorno de desarrollo.

\subsection{Base de Datos --- PostgreSQL y Docker Compose}

Se implementaron tres servicios en Docker Compose:

\begin{itemize}[leftmargin=*]
    \item \textbf{PostgreSQL 16:} Puerto 5433 (evitando conflicto con instalación
          del sistema en 5432), con usuario \texttt{vet} y base de datos
          \texttt{vetrinaria}.
    \item \textbf{MinIO:} Almacenamiento de objetos S3-compatible para archivos
          (imágenes, radiografías, facturas).
    \item \textbf{Redis:} Caché y manejo de sesiones.
\end{itemize}

Todos los contenedores se inician con \texttt{docker compose up -d} y comparten
una red Docker interna.

\subsection{Esquema de Base de Datos --- Drizzle ORM}

Se definieron 7 tablas usando \textbf{Drizzle ORM} con PostgreSQL:

\begin{table}[h]
\centering
\caption{Entidades del esquema de base de datos}
\begin{tabularx}{\textwidth}{|l|X|l|}
\hline
\textbf{Tabla} & \textbf{Propósito} & \textbf{Campos clave} \\
\hline
\texttt{users} & Usuarios del sistema & email, password\_hash, role \\
\texttt{clients} & Clientes (dueños de mascotas) & name, email, phone \\
\texttt{patients} & Pacientes (mascotas) & name, species, breed, client\_id \\
\texttt{appointments} & Citas veterinarias & date, status, patient\_id, user\_id \\
\texttt{medical\_records} & Historial clínico & diagnosis, treatment, patient\_id \\
\texttt{products} & Productos e inventario & name, price, stock \\
\texttt{suppliers} & Proveedores & name, contact, phone \\
\texttt{invoices} & Facturas & invoice\_number, total, status \\
\texttt{invoice\_items} & Partidas de factura & description, quantity, unit\_price \\
\texttt{notifications} & Notificaciones & title, message, type, read \\
\hline
\end{tabularx}
\end{table}

Las relaciones clave incluyen:

\begin{itemize}[leftmargin=*]
    \item \texttt{patients} → \texttt{clients} (FK: \texttt{client\_id})
    \item \texttt{appointments} → \texttt{patients}, \texttt{clients}, \texttt{users}
    \item \texttt{medical\_records} → \texttt{patients}, \texttt{appointments}, \texttt{users}
\end{itemize}

\subsection{TailwindCSS v4 y shadcn/ui}

Se implementó \textbf{TailwindCSS v4} con \textbf{shadcn/ui} para los componentes
de interfaz. Los componentes base creados son:

\begin{itemize}[leftmargin=*]
    \item \textbf{Button:} Variantes (default, destructive, outline, secondary,
          ghost, link) y tamaños (default, sm, lg, icon).
    \item \textbf{Card:} Contenedor con header, content y footer.
    \item \textbf{Input:} Campo de texto con validación.
    \item \textbf{Label:} Etiqueta accesible para formularios.
    \item \textbf{Avatar:} Imagen circular con fallback de iniciales.
    \item \textbf{Separator:} Línea divisoria horizontal.
    \item \textbf{DropdownMenu:} Menú contextual desplegable.
    \item \textbf{Textarea:} Área de texto multilínea.
\end{itemize}

Todos los componentes siguen la convención \texttt{cn()} para fusión de clases
Tailwind, usando \texttt{clsx} y \texttt{tailwind-merge}.

\subsection{Layout Responsivo}

Se implementó un layout de dos columnas:

\begin{itemize}[leftmargin=*]
    \item \textbf{Sidebar:} Colapsable en móviles, con navegación principal.
          Utiliza estado local para el toggle con icono \texttt{ChevronLeft}.
          Incluye enlaces a Dashboard, Agenda, Clientes, Pacientes, Historial
          Clínico, Vacunación, Facturación, Inventario y Mensajes.
    \item \textbf{Navbar:} Barra superior con notificaciones, cambio de tema
          (claro/oscuro con \texttt{next-themes}) y menú de usuario con avatar.
    \item \textbf{ThemeProvider:} Implementado con \texttt{next-themes}, soporta
          tema claro, oscuro y seguimiento del sistema.
\end{itemize}

El layout usa \textbf{route groups} de Next.js para separar la interfaz del
dashboard (ruta \texttt{(dashboard)} con sidebar/navbar) de las páginas de
autenticación (ruta \texttt{auth}, sin sidebar).

\subsection{CI/CD --- GitHub Actions}

Se configuró un pipeline de CI en \texttt{.github/workflows/ci.yml} que ejecuta
en cada push a \texttt{main} y cada Pull Request:

\begin{enumerate}[leftmargin=*]
    \item \textbf{Lint:} Verificación de código con ESLint.
    \item \textbf{Typecheck:} Validación de tipos con TypeScript.
    \item \textbf{Build:} Compilación de producción.
\end{enumerate}

\section{Sprint 1 --- Autenticación y Roles}

\subsection{NextAuth v5}

Se implementó \textbf{NextAuth v5} (beta 31) con el proveedor de credenciales
(email y contraseña) y estrategia de sesión JWT:

\begin{lstlisting}[language=TypeScript, caption=Configuración de NextAuth]
export const { handlers, signIn, signOut, auth } = NextAuth({
  providers: [
    Credentials({
      async authorize(credentials) {
        const user = await db.select().from(schema.users)
          .where(eq(schema.users.email, email))
          .then((rows) => rows[0]);

        const isValid = await compare(password, user.passwordHash);
        if (!isValid) return null;

        return { id: String(user.id), email: user.email,
                 name: ..., role: user.role };
      },
    }),
  ],
  session: { strategy: 'jwt' },
});
\end{lstlisting}

Los tipos de TypeScript se extendieron para incluir \texttt{role} e \texttt{id}
en los objetos \texttt{User}, \texttt{Session} y \texttt{JWT}.

\subsection{Páginas de Autenticación}

\subsubsection{Login (\texttt{/auth/login})}
Formulario \texttt{useState} con campos de email y contraseña. Utiliza
\texttt{signIn('credentials')} del lado del cliente. Muestra errores genéricos
``Credenciales inválidas'' sin revelar si el usuario existe. Redirige al
dashboard tras autenticación exitosa.

\subsubsection{Registro (\texttt{/auth/register})}
Formulario con nombre, apellido, email y contraseña. Envía los datos a
\texttt{POST /api/auth/register} que hashea la contraseña con bcrypt
(12 rondas) y crea el usuario con rol \texttt{receptionist} por defecto.

\subsection{Middleware de Protección}

Implementado en \texttt{middleware.ts} con el siguiente comportamiento:

\begin{itemize}[leftmargin=*]
    \item \textbf{Rutas públicas:} \texttt{/auth/login}, \texttt{/auth/register}
          --- accesibles sin autenticación.
    \item \textbf{Sin sesión:} Redirige a \texttt{/auth/login}.
    \item \textbf{Con sesión en ruta pública:} Redirige a \texttt{/dashboard}.
    \item \textbf{Protección por rol:} Clasificación de rutas por rol
          (admin/veterinarian tienen acceso completo; technician/receptionist
          tienen acceso restringido a módulos específicos).
\end{itemize}

\subsection{Integración con el Layout}

La sesión se integra mediante:

\begin{itemize}[leftmargin=*]
    \item \textbf{SessionProvider:} Envuelve toda la aplicación en el root layout.
    \item \textbf{Navbar:} Muestra nombre e iniciales del usuario autenticado,
          con menú desplegable para cerrar sesión.
    \item \textbf{Dashboard:} Verifica la sesión en el servidor con \texttt{auth()}
          y redirige a login si no está autenticado.
\end{itemize}

\subsection{Manejo de Errores}

El registro valida: campos obligatorios, email único, y errores internos
(HTTP 500). El login usa mensajes genéricos por seguridad.

\section{Sprint 2 --- Gestión de Clientes}

\subsection{Arquitectura}

El módulo de clientes sigue el patrón \textbf{Server Actions} de Next.js 15,
con un archivo de acciones en \texttt{lib/actions/clients.ts} y páginas
organizadas bajo la ruta \texttt{/clients}.

\subsection{Server Actions}

\begin{lstlisting}[language=TypeScript, caption=API de server actions para clientes]
export async function getClients(search?: string)
export async function getClient(id: number)
export async function createClient(data: ClientInput)
export async function updateClient(id: number, data: ClientInput)
export async function deleteClient(id: number)
\end{lstlisting}

Las operaciones utilizan Drizzle ORM con validación Zod. Las acciones
\texttt{createClient} y \texttt{updateClient} revalidan las rutas
\texttt{/clients} mediante \texttt{revalidatePath}.

\subsection{Páginas}

\begin{itemize}[leftmargin=*]
    \item \texttt{/clients} --- Lista con tabla responsiva, búsqueda por
          query string (\texttt{?q=}) sobre nombre, email y teléfono.
    \item \texttt{/clients/new} --- Formulario de creación con campos:
          nombre, email, teléfono, dirección, contacto de emergencia y notas.
    \item \texttt{/clients/[id]} --- Vista detalle con información de
          contacto, sección de mascotas (placeholder) y últimas citas.
    \item \texttt{/clients/[id]/edit} --- Formulario de edición precargado.
\end{itemize}

\subsection{Componente Reutilizable}

El \texttt{ClientForm} es un componente cliente (\texttt{'use client'}) que
recibe \texttt{defaultValues} y \texttt{onSubmit} como props. Utiliza estado
local con \texttt{useState} para manejo de carga y errores. Los campos se
procesan con \texttt{FormData} nativo del navegador.

\section{Sprint 3 --- Gestión de Pacientes}

\subsection{Esquema de Datos}

La tabla \texttt{patients} incluye:

\begin{itemize}[leftmargin=*]
    \item \texttt{client\_id} (FK $\to$ clients) --- Relación con el dueño.
    \item \texttt{name}, \texttt{species} (enum: dog/cat/rabbit/bird/reptile/other),
          \texttt{breed}, \texttt{gender} (enum: male/female).
    \item \texttt{birth\_date}, \texttt{color}, \texttt{microchip},
          \texttt{allergies}, \texttt{notes}.
    \item \texttt{weight} (JSONB) --- Historial de pesos.
    \item \texttt{photo\_url}, \texttt{active}, timestamps.
\end{itemize}

La relación con \texttt{clients} se define mediante \texttt{relations} de Drizzle.

\subsection{Server Actions}

\begin{lstlisting}[language=TypeScript, caption=API de server actions para pacientes]
export async function getPatients()
export async function getPatient(id: number)
export async function createPatient(data: PatientInput)
export async function updatePatient(id: number, data: PatientInput)
export async function deletePatient(id: number)
\end{lstlisting}

\subsection{Páginas}

\begin{itemize}[leftmargin=*]
    \item \texttt{/patients} --- Lista en formato tarjetas (\texttt{Card}) con
          grid responsivo, mostrando nombre, especie, raza, sexo y dueño.
    \item \texttt{/patients/new} --- Formulario con selector de dueño (carga
          desde \texttt{GET /api/clients/list}).
    \item \texttt{/patients/[id]} --- Vista detalle con información general
          (especie, raza, sexo, edad, color, microchip, alergias) y datos
          del dueño. Muestra notas si existen.
    \item \texttt{/patients/[id]/edit} --- Formulario de edición precargado.
\end{itemize}

\subsection{Componente Reutilizable}

El \texttt{PatientForm} utiliza un selector de dueños que se carga mediante
\texttt{fetch('/api/clients/list')} en un \texttt{useEffect}. Los campos de
especie y sexo usan \texttt{<select>} nativos con etiquetas en español.

\section{Sprint 4 --- Agenda de Citas}

\subsection{Esquema de Datos}

La tabla \texttt{appointments} incluye:

\begin{itemize}[leftmargin=*]
    \item \texttt{patient\_id} (FK $\to$ patients), \texttt{client\_id} (FK $\to$ clients),
          \texttt{vet\_id} (FK $\to$ users) --- Relaciones con paciente, dueño y veterinario.
    \item \texttt{date}, \texttt{start\_time}, \texttt{end\_time} --- Fecha y hora de la cita.
    \item \texttt{type} (enum: consultation/vaccination/surgery/...) --- Tipo de consulta.
    \item \texttt{status} (enum: scheduled/confirmed/checked\_in/...) --- Estado de la cita.
    \item \texttt{room}, \texttt{notes} --- Sala asignada y notas.
\end{itemize}

\subsection{Server Actions}

\begin{lstlisting}[language=TypeScript, caption=API de server actions para citas]
export async function getAppointments(date?: string)
export async function getAppointment(id: number)
export async function createAppointment(data: AppointmentInput)
export async function updateAppointment(id: number, data: AppointmentInput)
export async function updateAppointmentStatus(id: number, status: string)
\end{lstlisting}

Además del CRUD estándar, se implementó \texttt{updateAppointmentStatus}
para transiciones rápidas de estado directamente desde la lista y el detalle.

\subsection{Páginas}

\begin{itemize}[leftmargin=*]
    \item \texttt{/appointments?date=YYYY-MM-DD} --- Vista diaria con
          navegación entre fechas (botones {@literal <} Hoy {@literal >}),
          botón rápido ``Hoy'', y atajos de cambio de estado (2 acciones
          rápidas por cita). Muestra hora, paciente, dueño, tipo, sala y
          veterinario.
    \item \texttt{/appointments/new} --- Formulario con selectores
          dependientes: al elegir un dueño se cargan sus mascotas
          mediante \texttt{GET /api/patients/by-client/[id]}.
    \item \texttt{/appointments/[id]} --- Detalle completo con estado
          actual, datos de fecha/hora, responsables, paciente y notas.
          Incluye panel de transición de estado con los siguientes pasos
          válidos según el estado actual (por ejemplo: Programada
          $\to$ Confirmada $\to$ Check-In $\to$ En consulta $\to$ Check-Out).
    \item \texttt{/appointments/[id]/edit} --- Formulario de edición precargado.
\end{itemize}

\subsection{Transiciones de Estado}

El flujo de estados sigue esta máquina de estados simple:

\begin{verbatim}
scheduled --> confirmed --> checked_in --> in_exam --> checked_out
    |             |
    +--> cancelled+--> cancelled
\end{verbatim}

Cada estado solo permite transiciones hacia adelante, evitando cambios
inválidos. Los estados \texttt{cancelled}, \texttt{checked\_out} y
\texttt{no\_show} son terminales.

\subsection{Componente Reutilizable}

El \texttt{AppointmentForm} integra tres selectores que se cargan mediante
APIs:

\begin{itemize}[leftmargin=*]
    \item \texttt{GET /api/clients/list} --- Lista de dueños.
    \item \texttt{GET /api/patients/by-client/[id]} --- Mascotas del dueño
          seleccionado (carga dinámica con \texttt{useEffect}).
    \item \texttt{GET /api/users/vets} --- Veterinarios activos.
\end{itemize}

\section{Sprint 5 --- Historial Clínico}

\subsection{Esquema de Datos}

La tabla \texttt{medical\_records} sigue el formato \textbf{SOAP} (Subjetivo,
Objetivo, Evaluación, Plan):

\begin{itemize}[leftmargin=*]
    \item \texttt{patient\_id} (FK $\to$ patients), \texttt{appointment\_id}
          (FK $\to$ appointments), \texttt{created\_by} (FK $\to$ users).
    \item \texttt{subjective} (S) --- Historia y síntomas reportados.
    \item \texttt{objective} (O) --- Hallazgos del examen físico.
    \item \texttt{assessment} (A) --- Análisis y diagnóstico diferencial.
    \item \texttt{plan} (P) --- Tratamiento y seguimiento.
    \item \texttt{diagnosis} --- Diagnóstico principal.
    \item \texttt{vitals} (JSONB) --- Constantes vitales: peso, temperatura,
          frecuencia cardíaca y respiratoria.
\end{itemize}

\subsection{Server Actions}

\begin{lstlisting}[language=TypeScript, caption=API de server actions para historial]
export async function getMedicalRecords(patientId: number)
export async function getMedicalRecord(id: number)
export async function createMedicalRecord(data: MedicalRecordInput)
export async function updateMedicalRecord(id: number, data: MedicalRecordInput)
\end{lstlisting}

\subsection{Páginas}

\begin{itemize}[leftmargin=*]
    \item \texttt{/medical-records?patientId=X} --- Lista de registros de un
          paciente, ordenada del más reciente al más antiguo. Muestra
          diagnóstico, fecha, cita asociada y quién lo creó.
    \item \texttt{/medical-records/new?patientId=X\&appointmentId=Y} ---
          Formulario SOAP completo con campos para constantes vitales
          (peso, temperatura, FC, FR), selector de cita asociada y las
          cuatro secciones SOAP.
    \item \texttt{/medical-records/[id]} --- Vista detalle con diagnóstico
          destacado (borde izquierdo color primario), tarjetas de signos
          vitales en grid 4 columnas, y las cuatro secciones SOAP en grid
          2 columnas.
    \item \texttt{/medical-records/[id]/edit} --- Formulario de edición
          precargado.
\end{itemize}

\subsection{Integración con Pacientes}

La página de detalle del paciente (\texttt{/patients/[id]}) incluye ahora un
botón ``Historial'' que enlaza a
\texttt{/medical-records?patientId=X}, permitiendo acceso directo al
historial clínico desde el perfil del paciente.

\subsection{Componente Reutilizable}

El \texttt{MedicalRecordForm} implementa:

\begin{itemize}[leftmargin=*]
    \item Selector de cita asociada, cargado dinámicamente mediante
          \texttt{GET /api/appointments/by-patient/[id]}.
    \item Cuatro campos numéricos para constantes vitales en grid.
    \item Campo de diagnóstico como texto simple.
    \item Cuatro áreas de texto grandes para las secciones SOAP en grid
          2x2.
\end{itemize}

\section{Sprint 6 --- Facturación}

\subsection{Esquema de Datos}

Se agregaron las tablas \texttt{products}, \texttt{suppliers}, \texttt{invoices} e \texttt{invoice\_items}:

\begin{itemize}[leftmargin=*]
    \item \texttt{products} --- Productos y servicios con nombre, tipo (medication/supply/food/service),
          SKU, precio, precio de costo, cantidad en stock y punto de reorden.
    \item \texttt{suppliers} --- Proveedores con nombre, contacto, email, teléfono y dirección.
    \item \texttt{invoices} --- Facturas con folio único (\texttt{INV-XXXXXX}), subtotal,
          IVA (16\%), total, estado (draft/issued/paid/cancelled) y notas.
    \item \texttt{invoice\_items} --- Partidas de la factura con descripción, cantidad,
          precio unitario y total calculado.
\end{itemize}

\subsection{Server Actions}

Se implementaron tres conjuntos de acciones:

\begin{lstlisting}[language=TypeScript, caption=API de server actions para facturación]
// Productos
export async function getProducts()
export async function getProduct(id: number)
export async function createProduct(data: ProductInput)
export async function updateProduct(id: number, data: ProductInput)

// Proveedores
export async function getSuppliers()
export async function getSupplier(id: number)
export async function createSupplier(data: SupplierInput)
export async function updateSupplier(id: number, data: SupplierInput)

// Facturas
export async function getInvoices()
export async function getInvoice(id: number)
export async function createInvoice(data: InvoiceInput)
export async function updateInvoiceStatus(id: number, status: string)
\end{lstlisting}

La generación de números de factura es secuencial (\texttt{INV-000001},
\texttt{INV-000002}, \dots), consultando el último folio en la base de datos
e incrementando.

\subsection{Páginas}

\begin{itemize}[leftmargin=*]
    \item \texttt{/inventory} --- Lista de productos con tabla responsiva, alerta
          visual cuando el stock está por debajo del punto de reorden.
    \item \texttt{/inventory/new}, \texttt{/inventory/[id]/edit} --- Formulario
          de producto con campos: nombre, tipo (select), SKU, precio, precio de costo,
          stock y punto de reorden.
    \item \texttt{/suppliers} --- Lista en tarjetas con grid responsivo.
    \item \texttt{/suppliers/new}, \texttt{/suppliers/[id]/edit} --- Formulario
          de proveedor.
    \item \texttt{/invoices} --- Lista de facturas con tabla, colores por estado,
          botones de acción rápida (Emitir, Pagar).
    \item \texttt{/invoices/new} --- Formulario de factura con selección de
          cliente, cita opcional, conceptos dinámicos (agregar/eliminar filas),
          cálculo automático de subtotal, IVA y total.
    \item \texttt{/invoices/[id]} --- Detalle de factura con acciones según
          estado (Emitir, Pagar, Cancelar).
\end{itemize}

\subsection{Características Destacadas}

\begin{itemize}[leftmargin=*]
    \item \textbf{IVA fijo del 16\%} sobre el subtotal, aplicado automáticamente.
    \item \textbf{Transiciones de estado:} draft $\to$ issued $\to$ paid |
          cancelled, siguiendo el ciclo de vida real de una factura.
    \item \textbf{Conceptos dinámicos:} El formulario de factura permite agregar
          y eliminar filas de conceptos con cálculo en tiempo real.
    \item \textbf{Selectores dependientes:} Al seleccionar un cliente en la
          factura, se cargan sus citas disponibles como opción de asociación.
\end{itemize}

\section{Sprint 7 --- Dashboard con Métricas Reales}

\subsection{Server Action}

Se creó \texttt{lib/actions/dashboard.ts} con la función
\texttt{getDashboardMetrics()} que ejecuta las siguientes consultas en
paralelo mediante \texttt{Promise.all}:

\begin{itemize}[leftmargin=*]
    \item \textbf{Citas hoy:} Conteo total y agrupado por estado.
    \item \textbf{Total clientes:} Conteo simple.
    \item \textbf{Total pacientes:} Conteo simple.
    \item \textbf{Ingresos hoy:} Suma de facturas pagadas con fecha de creación
          hoy, usando \texttt{coalesce(sum(total),'0')}.
    \item \textbf{Ingresos del mes:} Suma de facturas pagadas desde el primer
          día del mes.
    \item \textbf{Próximas citas:} Join con pacientes y clientes, ordenado por
          fecha y hora, límite 5.
    \item \textbf{Actividad reciente:} Últimos 5 registros de clientes, pacientes
          y facturas, combinados y ordenados por fecha descendente (máximo 10).
\end{itemize}

\subsection{Páginas Actualizadas}

El dashboard (\texttt{/}) ahora muestra:

\begin{itemize}[leftmargin=*]
    \item \textbf{4 tarjetas de resumen} con datos en vivo: Citas Hoy (con
          desglose de pendientes), Clientes totales, Pacientes totales,
          Ingresos de hoy y del mes.
    \item \textbf{Panel de próximas citas} con cards enlazables mostrando hora,
          paciente, dueño, tipo y estado (con colores semánticos).
    \item \textbf{Panel de actividad reciente} con eventos de creación de
          clientes, pacientes y facturas, cada uno con identificador de color
          por tipo.
\end{itemize}

\subsection{Manejo de Estados Vacíos}

Cada sección del dashboard maneja correctamente el estado vacío con mensajes
apropiados y enlaces de acción.

\section{Sprint 8 --- Notificaciones}

\subsection{Esquema de Datos}

Se agregó la tabla \texttt{notifications}:

\begin{itemize}[leftmargin=*]
    \item \texttt{user\_id} (FK $\to$ users) --- Usuario destinatario.
    \item \texttt{title} --- Título de la notificación.
    \item \texttt{message} --- Cuerpo del mensaje.
    \item \texttt{type} (enum: appointment/invoice/client/patient/system) ---
          Categoría para el icono y filtrado.
    \item \texttt{link} --- URL opcional a la que redirige al hacer clic.
    \item \texttt{read} (boolean) --- Estado de lectura.
    \item \texttt{created\_at} --- Marca de tiempo.
\end{itemize}

Se generó la migración \texttt{0002\_notifications.sql} y se aplicó a la base
de datos.

\subsection{Server Actions y API}

\begin{itemize}[leftmargin=*]
    \item \texttt{createNotification(data)} --- Inserta una notificación para un
          usuario específico.
    \item \texttt{getNotifications()} --- Recupera las últimas 50 notificaciones
          del usuario autenticado.
    \item \texttt{getUnreadCount()} --- Conteo de notificaciones no leídas.
    \item \texttt{markAsRead(id)} --- Marca una notificación como leída.
    \item \texttt{markAllAsRead()} --- Marca todas como leídas.
    \item \texttt{notifyNewAppointment(...)} --- Crea notificaciones para el
          veterinario asignado y todos los administradores al agendar una cita.
    \item \texttt{notifyNewInvoice(...)} --- Crea notificaciones para
          administradores al crear una factura.
    \item \texttt{GET /api/notifications} --- Endpoint JSON para la Navbar
          (cliente), devuelve conteo y últimas 5 notificaciones.
\end{itemize}

\subsection{Navbar con Campana de Notificaciones}

La Navbar se actualizó para incluir:

\begin{itemize}[leftmargin=*]
    \item \textbf{Campana} con contador de no leídos (badge rojo con número).
          Muestra ``9+'' cuando hay más de 9.
    \item \textbf{Dropdown} con las últimas 5 notificaciones, cada una con
          icono según el tipo, título, mensaje truncado y punto azul si no
          está leída.
    \item \textbf{Actualización automática} cada 15 segundos mediante
          \texttt{setInterval} y \texttt{fetch} al endpoint.
    \item \textbf{Link ``Ver todas''} que redirige a la página de notificaciones.
    \item Al hacer clic en una notificación no leída, se marca automáticamente
          como leída mediante \texttt{markAsRead}.
\end{itemize}

\subsection{Página de Notificaciones}

\begin{itemize}[leftmargin=*]
    \item \texttt{/notifications} --- Página completa del lado del servidor con
          lista de todas las notificaciones.
    \item Muestra icono por tipo, título, mensaje completo, fecha formateada
          en español.
    \item Las no leídas tienen resaltado visual (borde y fondo tenue).
    \item Botón ``Ver'' si la notificación tiene enlace asociado.
    \item Botón ``Leído'' individual y ``Marcar todas leídas'' en el header.
    \item Estado vacío con icono y mensaje.
\end{itemize}

\subsection{Disparadores Automáticos}

Se integraron notificaciones en las siguientes acciones existentes:

\begin{itemize}[leftmargin=*]
    \item \texttt{createAppointment}: Notifica al veterinario asignado y a todos
          los administradores cuando se agenda una cita.
    \item \texttt{createClient}: Notifica a todos los administradores cuando se
          registra un nuevo cliente.
    \item \texttt{createInvoice}: Notifica a administradores y recepcionistas
          cuando se crea una factura.
\end{itemize}

\subsection{Navegación}

Se agregó el enlace ``Notificaciones'' en la Sidebar con el icono de campana
(\texttt{Bell} de lucide-react).

\section{Sprint 9 --- Backups y Restauración}

\subsection{Descripción General}

El sistema de backups permite crear, listar, restaurar y eliminar copias de
seguridad de la base de datos PostgreSQL directamente desde la interfaz web.
Los backups se generan en formato \texttt{custom} (\texttt{-Fc}) de pg\_dump,
que ofrece compresión y la capacidad de restaurar con procesos paralelos.

\subsection{Backup Utility (\texttt{lib/backup.ts})}

Módulo del lado del servidor (\texttt{'server-only'}) que ejecuta comandos
del sistema mediante \texttt{child\_process.exec}:

\begin{itemize}[leftmargin=*]
    \item \texttt{createBackup()} --- Ejecuta \texttt{pg\_dump -Fc} parseando
          la URL de conexión de \texttt{DATABASE\_URL}. Genera archivos con
          nombre \texttt{vetrinaria-YYYY-MM-DDTHH-MM-SS.dump} en el directorio
          \texttt{backups/}.
    \item \texttt{listBackups()} --- Lee el directorio \texttt{backups/},
          filtrando archivos \texttt{*.dump}, y devuelve metadatos (nombre,
          tamaño, fecha de modificación) ordenados del más reciente al más
          antiguo.
    \item \texttt{restoreBackup(filename)} --- Ejecuta \texttt{pg\_restore -c}
          (limpia objetos existentes antes de restaurar) usando el archivo
          seleccionado. Timeout de 10 minutos.
    \item \texttt{deleteBackup(filename)} --- Elimina un archivo de backup.
    \item \texttt{formatSize()} --- Convierte bytes a formato legible (B, KB, MB).
\end{itemize}

\subsection{Server Actions (\texttt{lib/actions/backups.ts})}

\begin{itemize}[leftmargin=*]
    \item \texttt{triggerBackup()} --- Crea un nuevo backup y revalida la ruta.
    \item \texttt{getBackups()} --- Lista los backups disponibles.
    \item \texttt{removeBackup(filename)} --- Elimina un backup.
    \item \texttt{restoreFromBackup(filename)} --- Restaura la base de datos
          desde el archivo seleccionado.
\end{itemize}

\subsection{Páginas}

\texttt{/settings} — Centro de configuración con tarjetas enlazables a
diferentes secciones (Backups, General, Seguridad, Apariencia).

\texttt{/settings/backups} — Gestión completa de backups:

\begin{itemize}[leftmargin=*]
    \item Botón \textbf{Crear Backup} que ejecuta pg\_dump y muestra el resultado.
    \item Lista de backups con nombre, tamaño, fecha y acciones:
          \textbf{Restaurar} (con advertencia de sobreescritura) y
          \textbf{Eliminar}.
    \item Panel de advertencia sobre la restauración (card ámbar).
    \item Panel de configiración de \textbf{Backups Automáticos} con
          instrucciones para programar via cron.
    \item Estado vacío con mensaje y llamado a acción.
\end{itemize}

\subsection{Backups Automáticos (Cron)}

Se incluye el script \texttt{scripts/backup.sh} que puede ejecutarse
via cron para backups diarios automatizados:

\begin{lstlisting}[language=bash, caption=Programar backup diario a las 2:00 AM]
0 2 * * * /path/to/vetrinaria/scripts/backup.sh
\end{lstlisting}

El script parsea \texttt{DATABASE\_URL} de \texttt{.env}, ejecuta pg\_dump
con formato personalizado y almacena el archivo en \texttt{backups/}.

\subsection{Seguridad}

\begin{itemize}[leftmargin=*]
    \item El módulo de backup está marcado como \texttt{'server-only'},
          impidiendo su importación desde el cliente.
    \item Las server actions verifican la sesión mediante \texttt{auth()}
          (el middleware de NextAuth protege todas las rutas del dashboard).
    \item La restauración muestra una advertencia explícita antes de
          proceder, ya que sobreescribe todos los datos actuales.
\end{itemize}

\section{Sprint 10 --- Multi-Tenant}

\subsection{Arquitectura}

Se implementó una arquitectura multi-tenant basada en filas (row-based), donde
cada clínica (tenant) comparte las mismas tablas pero sus datos se aíslan
mediante la columna \texttt{clinic\_id} en todas las entidades del sistema.

\subsection{Esquema de Datos}

La tabla \texttt{clinics} almacena la información de cada clínica:

\begin{itemize}[leftmargin=*]
    \item \texttt{name}, \texttt{slug} (único, para identificación), \texttt{email},
          \texttt{phone}, \texttt{address}.
    \item \texttt{settings} (JSONB) --- Configuración específica (tema, locale,
          timezone, prefijo de factura, tasa de impuesto).
    \item \texttt{active}, timestamps de creación y actualización.
\end{itemize}

Se agregó \texttt{clinic\_id} (FK $\to$ clinics) a todas las tablas de datos:
\texttt{users}, \texttt{clients}, \texttt{patients}, \texttt{appointments},
\texttt{medical\_records}, \texttt{products}, \texttt{suppliers},
\texttt{invoices}.

\subsection{Migración}

La migración \texttt{0003\_multi-tenant.sql} maneja datos existentes:

\begin{enumerate}[leftmargin=*]
    \item Crea la tabla \texttt{clinics} e inserta una clínica por defecto
          (\texttt{Clínica Principal}).
    \item Agrega \texttt{clinic\_id} como nullable en todas las tablas.
    \item Actualiza los registros existentes con el ID de la clínica por defecto.
    \item Establece la restricción \texttt{NOT NULL}.
    \item Agrega las llaves foráneas correspondientes.
\end{enumerate}

\subsection{Rol Super Admin}

Se agregó el rol \texttt{super\_admin} al enum de roles de usuario. Este rol
tiene acceso completo a todas las clínicas y a la administración del sistema,
incluyendo la gestión de clínicas.

\subsection{Clinic Context Utility (\texttt{lib/clinic.ts})}

\begin{itemize}[leftmargin=*]
    \item \texttt{getCurrentClinicId()} --- Obtiene el ID de la clínica desde
          la sesión JWT del usuario autenticado.
    \item \texttt{getCurrentUserRole()} --- Obtiene el rol del usuario actual.
    \item \texttt{isSuperAdmin()} --- Verifica si el usuario es super admin.
\end{itemize}

Todas las funciones están marcadas como \texttt{'server-only'}.

\subsection{Autenticación Actualizada}

El archivo \texttt{lib/auth.ts} se actualizó para incluir \texttt{clinicId}
en el JWT y la sesión. Los tipos de TypeScript se extendieron para incluir
\texttt{clinicId} en \texttt{User} y \texttt{Session}.

El middleware se actualizó para reconocer el rol \texttt{super\_admin}.

\subsection{Server Actions (\texttt{lib/actions/clinics.ts})}

\begin{itemize}[leftmargin=*]
    \item \texttt{getClinics()} --- Lista todas las clínicas.
    \item \texttt{getClinic(id)} --- Obtiene una clínica por ID.
    \item \texttt{createClinic(data)} --- Crea una nueva clínica (con validación
          Zod del slug).
    \item \texttt{updateClinic(id, data)} --- Actualiza una clínica existente.
\end{itemize}

\subsection{Páginas}

\texttt{/settings/clinics} — Lista de clínicas registradas:

\begin{itemize}[leftmargin=*]
    \item Cada clínica muestra nombre, slug y email.
    \item Botón para editar cada clínica.
    \item Botón ``Nueva Clínica'' para crear.
    \item Estado vacío con mensaje.
\end{itemize}

\texttt{/settings/clinics/new} — Formulario de creación con campos: nombre,
slug (validado con regex \texttt{[a-z0-9-]+}), email, teléfono y dirección.

\texttt{/settings/clinics/[id]/edit} — Formulario de edición precargado.

\subsection{Componente Reutilizable}

El \texttt{ClinicForm} es un componente cliente con campo slug validado,
grid responsivo para email/teléfono, y manejo de errores.

\subsection{Filtrado por Clínica en Consultas}

Se actualizaron las consultas críticas para filtrar por \texttt{clinic\_id}:

\begin{itemize}[leftmargin=*]
    \item \textbf{Dashboard:} Todas las métricas (citas, clientes, pacientes,
          ingresos, actividad reciente) filtradas por la clínica del usuario.
    \item \textbf{Clientes:} Lista y creación filtradas/asignadas a la clínica.
    \item \textbf{Pacientes:} Lista y creación filtradas/asignadas a la clínica.
    \item \textbf{Citas:} Lista y creación filtradas/asignadas a la clínica.
\end{itemize}

\subsection{Datos Semilla}

El script \texttt{seed.ts} se actualizó para:

\begin{itemize}[leftmargin=*]
    \item Crear una clínica por defecto (\texttt{Clínica Veterinaria Central}).
    \item Asignar el usuario admin a esa clínica con rol \texttt{super\_admin}.
    \item Asignar clientes y pacientes de ejemplo a la clínica.
\end{itemize}

\section{Tecnologías Utilizadas}

\begin{table}[h]
\centering
\caption{Stack tecnológico del proyecto}
\begin{tabularx}{\textwidth}{|l|X|}
\hline
\textbf{Tecnología} & \textbf{Versión} \\
\hline
Next.js & 15.5.18 \\
TypeScript & 5.8 \\
TailwindCSS & 4 \\
pnpm & 11.5 \\
Turborepo & 2.9 \\
PostgreSQL & 16 \\
Drizzle ORM & 0.41 \\
NextAuth & 5.0.0-beta.31 \\
Docker Compose & 3.8 \\
Node.js & 22 \\
\hline
\end{tabularx}
\end{table}

Además: React 19, bcryptjs, clsx, tailwind-merge, lucide-react,
next-themes, zod.

\section{Estructura de Archivos --- Sprints 2 a 10}

\begin{lstlisting}[language=bash, caption=Archivos creados en Sprints 2 a 10]
apps/web/src/
+-- lib/actions/
|   +-- clients.ts           # CRUD clientes (5 server actions)
|   +-- clients-list.ts      # Lista simple para selects
|   +-- patients.ts          # CRUD pacientes (5 server actions)
|   +-- appointments.ts      # CRUD citas (5 server actions + status)
|   +-- medical-records.ts   # CRUD historial clinico (4 server actions)
|   +-- products.ts          # CRUD productos (4 server actions)
|   +-- suppliers.ts         # CRUD proveedores (4 server actions)
|   +-- invoices.ts          # CRUD facturas (4 server actions)
|   +-- dashboard.ts         # Metricas del dashboard (1 server action)
|   +-- notifications.ts     # Notificaciones (7 server actions)
|   +-- backups.ts           # Backups (4 server actions)
|   +-- clinics.ts           # CRUD clinicas (4 server actions)
+-- lib/
|   +-- backup.ts            # Utilidad de pg_dump (server-only)
|   +-- clinic.ts            # Contexto multi-tenant (server-only)
+-- components/
|   +-- clinic-form.tsx      # Formulario reutilizable de clinicas
+-- components/
|   +-- client-form.tsx      # Formulario reutilizable de clientes
|   +-- patient-form.tsx     # Formulario reutilizable de pacientes
|   +-- appointment-form.tsx # Formulario reutilizable de citas
|   +-- medical-record-form.tsx # Formulario SOAP con vitales
|   +-- product-form.tsx     # Formulario reutilizable de productos
|   +-- supplier-form.tsx    # Formulario reutilizable de proveedores
|   +-- ui/textarea.tsx      # Componente Textarea
+-- app/(dashboard)/
|   +-- clients/             # CRUD completo (list, new, [id], edit)
|   +-- patients/            # CRUD completo (list, new, [id], edit)
|   +-- appointments/        # CRUD completo (list, new, [id], edit)
|   +-- medical-records/     # CRUD completo (list, new, [id], edit)
|   +-- inventory/           # Lista, nuevo, editar
|   +-- suppliers/           # Lista, nuevo, editar
|   +-- invoices/            # Lista, nuevo, detalle
|   +-- notifications/       # Lista completa de notificaciones
|   +-- settings/
|   |   +-- page.tsx         # Centro de configuracion
|   |   +-- backups/
|   |       +-- page.tsx     # Gestion de backups
|   +-- page.tsx             # Dashboard con metricas reales
+-- app/api/
|   +-- clients/list/route.ts       # API lista de clientes
|   +-- patients/by-client/[id]/route.ts  # API mascotas por dueno
|   +-- appointments/by-patient/[id]/route.ts # API citas por paciente
|   +-- users/vets/route.ts        # API lista de veterinarios
|   +-- notifications/route.ts     # API notificaciones (GET)
+-- packages/db/src/schema/
|   +-- clinics.ts           # Esquema de clinicas (multi-tenant)
+-- app/(dashboard)/settings/
|   +-- page.tsx             # Centro de configuracion
|   +-- backups/
|       +-- page.tsx         # Gestion de backups
|   +-- clinics/
|       +-- page.tsx         # Lista de clinicas
|       +-- new/page.tsx     # Crear clinica
|       +-- [id]/edit/page.tsx  # Editar clinica
+-- scripts/
|   +-- backup.sh           # Script para backup via cron
\end{lstlisting}

\section{Próximos Pasos}

\begin{enumerate}[leftmargin=*]
    \item \textbf{Sprint 10:} Multi-tenant --- aislamiento de datos
          entre clínicas con schema-based o row-level security.
    \item \textbf{Sprint 11:} Despliegue en producción con Docker
          Compose, HTTPS, y monitoreo.
\end{enumerate}

\end{document}
